% =============================================================================
% ЛЕКЦИЯ 1
% =============================================================================
\lecturedate{\today}

\section{Математическая логика}

Булевы функции, это штуки которые принимают на вход набор из нулей и единиц, а на выходе выдают тоже ноль или единицу.
\begin{theorem}
  Всего различных булевых функций от n переменных - $2^{2^n}$
\end{theorem}
\begin{proof}
  Каждая функция определяется своей таблицей истинности, в которой $2^n$ строк (каждая строка - это набор из нулей и единиц длины n). \\
  В каждой строке функция может принимать значение 0 или 1, всего вариантов - $2^{2^n}$.
\end{proof}
\begin{theorem}{Теорема Поста}
Система функций назвается полной, если она не содержится не в одном из пяти замкнутых классов булевых функций: \\
1. Класс функций, сохраняющих 0: $f(0, 0, ..., 0) = 0$ \\
2. Класс функций, сохраняющих 1: $f(1, 1, ..., 1) = 1$ \\
3. Класс линейных функций: функции которые могут быть выражены через операции ИЛИ по модулю 2 и И(Полином Жегалкина). \\
4. Класс самоддвойственных функций $\forall \alpha \in$ набор из нулей и единиц: $f(\alpha) = ¬f(\neg \alpha)$\\
5. Класс монотонных функций: $\forall \alpha \beta \in$ набор булевых прееменных: если $\alpha \leq \beta$ (т.е. для каждой переменной из $\alpha$ значение меньше или равно значению из $\beta$), то $f(\alpha) \leq f(\beta)$ \\
\end{theorem}
\begin{definition}{Закон деморгана}
  \\
  $\overline{A \land B} = \overline{A} \lor \overline{B}$ \\
  $\overline{A \lor B} = \overline{A} \land \overline{B}$
\end{definition}
\begin{definition}
  Замкнутый класс булевых функций - это такой класс булевых функций, что при композиции функций из этого класса получается функция, которая тоже принадлежит этому классу. (типо группа получатеся, шок)
\end{definition}

\subsection{Норм. формы}
\begin{definition}{Дизъюнктивная нормальная форма} 
  \\
  ДНФ - это набор коннюнкций, объединённых дизъюнкцией. Ну еще можно отрицание от некоторых брать, легко запомнить по названию, главная штука всегда написана в названии \\
  Пример: $(x_1 \land x_2 \land \overline{x_3}) \lor (\overline{x_1} \land x_2 \land x_3)$
\end{definition}
\begin{definition}
  {Конъюнктивная нормальная форма} \\
  КНФ - это набор дизъюнкциий, которые связаны конъюнкцией. \\
  Пример: $(x_1 \lor \overline{x_2} \lor x_3) \land (x_1 \lor x_2 \lor \overline{x_3})$
\end{definition}
Они еще связаны между собой, если взять отрицание от КНФ, то получим ДНФ и наоборот. Через закон Де Моргана можно доказать. (читателю етсь шанс потренироваться)
\begin{definition}{Тавтология}
Тавтология - это формула, которая возвращает 1 не зависимо от значений переменных.\\
Пример: $A \lor \overline{A}$ \\
\end{definition}
\begin{definition}{выполнимая формула}
Это когда хотя бы раз истина будет. \\
Пример: $A \land B$ \\
\end{definition}
\begin{definition}{невыполнимая формула}
\\
$\neg$ выполнимая формула(никогда 1 не получим)
\end{definition}
По ДНФ мы можем найти хотя бы один набор переменных, при котором формула будет истинной, а по КНФ мы можем найти набор переменных, при которых формула будет ложной. \\

\subsection{Совершенные норм формулы}
\begin{definition}{Совершенная дизъюнктивная нормальная форма}
  \\
  СДНФ - это ДНФ, в которой каждая конъюнкция содержит все переменные формулы. \\
  Пример: $(x_1 \land x_2 \land \overline{x_3}) \lor (\overline{x_1} \land x_2 \land x_3) \lor (x_1 \land \overline{x_2} \land x_3)$
\end{definition}
\begin{definition}{Совершенная конъюнктивная нормальная форма}
  \\
  СКНФ - это КНФ , в которой каждая дизъюнкция содержит все переменные формулы. \\
  Пример: $(x_1 \lor \overline{x_2} \lor x_3) \land (x_1 \lor x_2 \lor \overline{x_3}) \land (\overline{x_1} \lor x_2 \lor x_3)$
\end{definition}

через СКНФ и СДНФ можно выразить любую булеву функцию (так мы доказываем что через И/ИЛИ/НЕ можно выразить любую булеву функцию) \\

\subsection*{Мега гайд как построить СДНФ по таблице истинности} 

Пусть есть таблица истинности формата:
\begin{table}[h!]
\centering
\begin{tabular}{|c|c|c|c|}
\hline
$x_1$ & $x_2$ & $f (x_1, x_2)$ \\
\hline
0 & 0 & 1 \\
\hline
0 & 1 & 0 \\
\hline
1 & 0 & 1 \\
\hline
1 & 1 & 0 \\
\hline
\end{tabular}
\end{table}
теперь берем все строки, в которых функция равна 1, и для каждой такой строки строим конъюнкцию из переменных, где если переменная равна 0, то берем ее отрицание, а если 1, то просто переменную. \\
в итоге у нас получится: \\
для первой строки: $\overline{x_1} \land \overline{x_2}$ \\
для третьей строки: $x_1 \land \overline{x_2}$ \\
теперь берем все эти конъюнкции и объединяем их дизъюнкцией: \\
итог: $(\overline{x_1} \land \overline{x_2}) \lor (x_1 \land \overline{x_2})$ \\  
А теперь заюзаем закон Де Моргана и получим СКНФ: \\
$\overline{(\overline{x_1} \land \overline{x_2})} \land \overline{(x_1 \land \overline{x_2})}$ \\
что равносильно \\
$(x_1 \lor x_2) \land (\overline{x_1} \lor x_2)$ \\

\subsection{Полином Жегалкина}
В этой штуке есть три функции: 1 (константа), ИЛИ по модулю 2 (обозначается как $\oplus$) и И (обозначается как $\land$). \\
\begin{definition}
  {Полином Жегалкина} \\
  Полином Жегалкина - это представление булевой функции в виде суммы по модулю 2 (т.е. ИЛИ по модулю 2) нескольких конъюнкций переменных и константы 1. \\
  Пример: $a_{000} \oplus (a_{100} \land x_1) \oplus (a_{010} \land x_2) \oplus (a_{110} \land x_1 \land x_2) \oplus (a_{111} \land x_1 \land x_2 \land x_3) \oplus (a_{101} \land x_1 \land x_3)$
  Где индекс у а - это набор из нулей и единиц, которые показывают какие переменные участвуют в конъюнкции. \\
  Например, $a_{101}$ - это конъюнкция из $x_1$ и $x_3$, потому что в индексе на первом и третьем месте стоят единицы. \\
\end{definition}
\subsection*{Гайдик как построить полином Жегалкина по таблице истинности}

\subsection*{Гайдик как построить по ДНФ(КНФ) полином Жегалкина}
ну очеивидно что кнф и днф можно шафлить между собой и потом уже юзать один из гайдов
Из кнф в полиномчик Жегалкина: \\
1. Берем первую дизъюнкцию из ДНФ \\
2. Заметим что $x \oplus 1 = \overline{x}$ \\ значит все отрицания формата $\overline{x_1} \lor \overline{x_2} \lor ... \lor \overline{x_n}$ можно записать как $(x_1 \oplus 1) \lor (x_2 \oplus 1) \lor ... \lor (x_n \oplus 1)$ \\
3. Раскроем скобки по закону Де Моргана \\
4. Получим набор конъюнкций, которые нужно будет сложить по модулю 2 \\
5. Повторяем для всех дизъюнкций из ДНФ \\
6. Складываем все полученные конъюнкции по модулю 2 \\
7. Готово \\
